\section{Chernoff-Hoeffding Bound}

We use the following concentration inequality.


\begin{theorem}[Real-vaued Additive Chernoff-Hoeffding
  Bound]\label{chernoff}
  Let $X_1, X_2, \ldots, X_m$ be i.i.d. random variables with
  $\Ex{}{X_i} = \mu$ and $a\leq X_i \leq b$ for all $i$. Then for
  every $\alpha  > 0$,
  \[
    \Pr\left[ \left| \frac{\sum_i X_i}{m} - \mu \right| \geq \alpha
    \right] \leq 2 \exp\left(\frac{-2\alpha^2 m }{(b-a)^2} \right)
  \]
\end{theorem}

\iffalse
\begin{theorem}[McDiarmid Inequality]\label{mcdiarmid}
  Let $X_1, X_2, \ldots, X_m$ be independent random variables in the
  set $S$. Let $f\colon S^m \rightarrow \RR$ be a function of
  $X_1, \ldots , X_m$ with the following property: for any $i$ and for
  any $x_1, \ldots, x_m, x_i' \in S$,
  \[
    |f(x_1, \ldots, x_i, \ldots, x_m) - f(x_1, \ldots, x_i', \ldots,
    x_m)| \leq c_i.
  \]
  Then for any $\eps > 0$,
  \[
    \Pr[f - \Ex{}{f} \geq \eps ] \exp\left(
      \frac{-2\eps^2}{\sum_{i=1}^m c_i^2}\right)
  \]
\end{theorem}


\begin{theorem}[Real-valued non-i.i.d.~Chernoff-Hoeffding
  Bound]\label{chernoff2}
  Let $X_1, X_2, \ldots, X_m$ be independent random variables raning
  in $[0, 1]$, and $\beta>0$. Let $\mu = \sum_{i\in [n]}\Ex{}{X_i}$
  and suppose that $\sum_{i\in [n]} \mathrm{Var}(X_i) \leq
  \beta$. Then for every $\alpha \in  (0, 2\beta]$, it holds that
  \[
    \Pr\left[ \left| \sum_{i\in [n]} X_i - \mu \right| > \alpha
    \right] < 2\exp\left(\frac{-\alpha^2}{4\mu}\right)
  \]% take beta = \mu
\end{theorem}
\fi

\section{Generalization Bounds}
\label{app:generalization}

\begin{proof}[Proof of Theorems \ref{thm:SP-uc} and \ref{thm:FP-uc}]
We give a proof of Theorem \ref{thm:SP-uc}. The proof of Theorem \ref{thm:FP-uc} is identical, as false positive rates are just positive classification rates on the subset of the data for which $y = 0$.

Given a set of classifiers $\cH$ and protected groups $\cG$, define the following function class:
$$\mathcal{F}_{\cH, \cG} = \{f_{h,g}(x) \doteq h(x) \wedge g(x) : h \in \cH, g \in \cG\}$$
We can relate the VC-dimension of $\mathcal{F}_{\cH,\cG}$ to the VC-dimension of $\cH$ and $\cG$:
\begin{claim}
$$\mathrm{VCDIM}(\mathcal{F}_{\cH, \cG}) \leq \tilde{O}(\mathrm{VCDIM}(\cH) + \mathrm{VCDIM}(\cG))$$
\end{claim}
\begin{proof}
Let $S$ be a set of size $m$ shattered by $\mathcal{F}_{\cH, \cG}$. Let $\pi_{\mathcal{F}_{\cH, \cG}}(S)$ be the number of labelings of $S$ realized by elements of $\mathcal{F}_{\cH, \cG}$. By the definition of shattering, $\pi_{\mathcal{F}_{\cH, \cG}}(S) = 2^{m}$. Now for each labeling of $S$ by an element in $\mathcal{F}_{\cH, \cG}$, it is realized as $(f \wedge g)(S)$ for some $f \in \mathcal{F}, g \in \mathcal{G}$. But $(f \wedge g)(S) = f(S) \wedge g(S)$, and so it can be realized as the conjunction of a labeling of $S$ by an element of $\mathcal{F}$ and an element of $\mathcal{G}$. But since there are $\pi_{\mathcal{F}}(S)\pi_{\mathcal{G}}(S)$ such pairs of labelings, this immediately implies that $\pi_{\mathcal{F}_{\cH, \cG}}(S) \leq \pi_{\mathcal{F}}(S)\pi_{\mathcal{G}}(S)$. Now by the Sauer-Shelah Lemma (see e.g. \cite{KV94}), $\pi_{\mathcal{F}}(S) = O(m^{\mathrm{VCDIM}(\cH)}), \pi_{\cG}(S) = O(m^{\mathrm{VCDIM}(\cG)})$. Thus $\pi_{\mathcal{F}_{\cH, \cG}}(S) = 2^{m} \leq O(m^{\mathrm{VCDIM}(\cH) + \mathrm{VCDIM}(\cG)})$, which implies that $m = \tilde{O}(\mathrm{VCDIM}(\cH) + \mathrm{VCDIM}(\cG))$, as desired.


\end{proof}
This bound, together with a standard VC-Dimension based uniform convergence theorem (see e.g. \cite{KV94}) implies that with probability $1-\delta$, for every $f_{h,g} \in \mathcal{F}_{\cH, \cG}$:
$$\left|\E_{(X,y) \sim \mathcal{P}}[f_{h,g}(X)] - \E_{(X,y) \sim \cP_S}[f_{h,g}(X)] \right| \leq \tilde{O}\left(\sqrt{\frac{(\mathrm{VCDIM}(\cH) + \mathrm{VCDIM}(\cG))\log m + \log(1/\delta)}{m}} \right)$$
Note that the left hand side of the above inequality can be written as:
$$\left|\Pr_{(X,y) \sim \mathcal{P}}[h(X) = 1 | g(x) = 1] \cdot \Pr_{(X,y) \sim \mathcal{P}}[g(x) = 1] -  \Pr_{(X,y) \sim \cP_S}[h(X) = 1 | g(x) = 1] \cdot \Pr_{(X,y) \sim \cP_S}[g(x) = 1] \right|$$
This completes our proof.
\iffalse
Thus, dividing both sides by  $\Pr_{(X,y) \sim \mathcal{P}}[g(x) = 1]$ and noting that by assumption this quantity is at least $\alpha$, we obtain:
$$\left| \Pr_{(X,y) \sim \mathcal{P}}[h(X) = 1 | g(x) = 1] - \Pr_{(X,y) \sim \cP_S}[h(X) = 1 | g(x) = 1]\cdot \frac{ \Pr_{(X,y) \sim \cP_S}[g(x) = 1]}{\Pr_{(X,y) \sim \mathcal{P}}[g(x) = 1]}\right| \leq $$
$$ \tilde{O}\left(\frac{1}{\alpha}\sqrt{\frac{(\mathrm{VCDIM}(\cH) + \mathrm{VCDIM}(\cG))\log m + \log(1/\delta)}{m}} \right)$$
By the triangle inequality, we can bound
$$\left| \Pr_{(X,y) \sim \mathcal{P}}[h(X) = 1 | g(x) = 1] - \Pr_{(X,y) \sim \cP_S}[h(X) = 1 | g(x) = 1]\right| \leq $$
$$\left| \Pr_{(X,y) \sim \mathcal{P}}[h(X) = 1 | g(x) = 1] - \Pr_{(X,y) \sim \cP_S}[h(X) = 1 | g(x) = 1]\cdot \frac{ \Pr_{(X,y) \sim \cP_S}[g(x) = 1]}{\Pr_{(X,y) \sim \mathcal{P}}[g(x) = 1]}\right| + $$
$$\left|\Pr_{(X,y) \sim \cP_S}[h(X) = 1 | g(x) = 1]\cdot \frac{ \Pr_{(X,y) \sim \cP_S}[g(x) = 1]}{\Pr_{(X,y) \sim \mathcal{P}}[g(x) = 1]} -  \Pr_{(X,y) \sim \cP_S}[h(X) = 1 | g(x) = 1]\right| $$
We have already bounded the first term: it remains to bound the second term.
Again, invoking a standard VC-dimension bound, we obtain:
$$\Pr_{(X,y) \sim \cP_S}[g(x) = 1] \leq \Pr_{(X,y) \sim \mathcal{P}}[g(x) = 1] +  O\left(\sqrt{\frac{\mathrm{VCDIM}(\cG)\log m + \log(1/\delta)}{m}} \right)$$
Since by assumption we have $\Pr_{(X,y) \sim \mathcal{P}}[g(x) = 1] \geq \alpha$, we therefore know that with probability $1-\delta$:
$$\frac{\Pr_{(X,y) \sim \cP_S}[g(x) = 1]}{\Pr_{(X,y) \sim \mathcal{P}}[g(x) = 1] } \leq 1+  O\left(\frac{1}{\alpha}\sqrt{\frac{\mathrm{VCDIM}(\cG)\log m + \log(1/\delta)}{m}} \right)$$
Finally, this lets us bound the 2nd term above:
$$\left|\Pr_{(X,y) \sim \cP_S}[h(X) = 1 | g(x) = 1]\cdot \frac{ \Pr_{(X,y) \sim \cP_S}[g(x) = 1]}{\Pr_{(X,y) \sim \mathcal{P}}[g(x) = 1]} -  \Pr_{(X,y) \sim \cP_S}[h(X) = 1 | g(x) = 1]\right| \leq O\left(\frac{1}{\alpha}\sqrt{\frac{\mathrm{VCDIM}(\cG)\log m + \log(1/\delta)}{m}} \right)$$
which completes the proof.
\fi
\end{proof}

\iffalse
\section{Relaxing the Base Rate Condition}
\sw{todo: need update to $\gamma$-formulation}

\begin{lemma}\label{dis-acc2}
  Suppose that the base rate
  $b_{SP}(D, P) \in [1/2 - \eps, 1/2 + \eps]$ and there exists a
  function $f$ such that $\Pr[f(x) = 1] = \alpha$ and
  $|\Pr[D(X) = 1 \mid f(x) = 1] - b_{SP}(D, P)| > \beta$.  Then
  $$
  \max\{\Pr[D(X) = f(x)], \Pr[D(X) = \neg f(x)]\} \geq
  (b_{SP} - 2\eps)  + \alpha\beta
  $$
\end{lemma}



\begin{proof}
  Let $b = b_{SP}(D, p)$ denote the base rate. Given that
  $|\Pr[D(X) = 1 \mid f(x) = 1] - b| > \beta$, we know either
  $\Pr[D(X) = 1 \mid f(x) = 1] > b+ \beta$ or
  $\Pr[D(X) = 1\mid f(x) = 1] < b - \beta$.

  In the first case, we know $\Pr[D(X) = 1\mid f(x) = 0] < b$, and
  so $\Pr[D(X) = 0 \mid f(x) = 0] > 1 - b$. It follows that
  \begin{align*}
    \Pr[D(X) = f(x) ] &= \Pr[D(X) = f(x) = 1] +\Pr[D(X) =
      f(x) = 0]\\
    &= \Pr[D(X) = 1 \mid f(x) = 1]\Pr[f(x) = 1] + \Pr[D(X) = 0 \mid f(x) = 0] \Pr[f(x)=0]\\
    &> \alpha (b + \beta) + (1-\alpha) (1 - b) =
      (\alpha - 1) b + (1 - \alpha) (1 - b) + b + \alpha\beta = (1 - \alpha)(1 - 2b) + b + \alpha\beta
  \end{align*}

  In the second case, we have
  $\Pr[D(X) = 0 \mid f(x) = 1] > (1 - b) + \beta$ and
  $\Pr[D(X) = 1\mid f(x) = 0] > b$. We can then bound
  \begin{align*}
   \Pr[D(X) = f(x) ] &= \Pr[D(X) = 1 \mid f(x) = 0]\Pr[f(x) = 0] + \Pr[D(X) = 0 \mid f(x) = 1] \Pr[f(x)=1]\\
    &> (1- \alpha) b  + \alpha (1 - b + \beta) =  \alpha(1 - 2b) +
      b+ \alpha\beta.
  \end{align*}

  In both cases, we have, and $(1 - 2b) \in [-2\eps, 2\eps]$ based our
  assumption on the base rate. Since $(1 - \alpha), \alpha\in [0,1]$,
  we have
  $$
  \max\{\Pr[D(X) = f(x)], \Pr[D(X) = \neg f(x)]\} \geq b - 2\eps +
  \alpha\beta
  $$
\end{proof}



\begin{lemma}\label{acc-dis2}
  Suppose that the base rate
  $b_{SP}(D, P) \in [1/2 - \eps, 1/2 + \eps]$ and there exists a
  function $f$ such that $\Pr[D(X) = f(x)] \geq b_1(D,P) + \gamma$ for
  some value $\gamma \in (0, 1/2)$. Then there exists a function $g$
  such that $\Pr[ g(x) = 1 ] \geq \gamma - 2\eps$ and
  $|\Pr[D(X) = 1 \mid g(x) = 1] - b_{SP}(D, p)| > \gamma - 2\eps$,
  where $g \in \{f, \neg f\}$.
\end{lemma}


\begin{proof}
  Let $b = b_{SP}(D, p)$ denote the base rate. Again recall that
  \begin{align*}
    \Pr[D(X) = f(x) ] &= \Pr[D(X) = f(x) = 1] +\Pr[D(X) =
                        f(x) = 0]\\
                      &= \Pr[D(X) = 1 \mid f(x) = 1]\Pr[f(x) = 1] + \Pr[D(X) = 0 \mid f(x) = 0] \Pr[f(x)=0]
  \end{align*}
  Since $\Pr[D(X) = f(x)] \geq b + \gamma$ and
  $\Pr[f(x) = 1] + \Pr[f(x) = 0] = 1$,
  \[
    \max\{ \Pr[D(X) = 1 \mid f(x) = 1], \Pr[D(X) = 0 \mid f(x)
    = 0]\} \geq b + \gamma.
  \]
  Thus, we must have either
\[
  \Pr[D(X) = 1 \mid f(x) = 1] \geq b+ \gamma \qquad \mbox{ or } \qquad
  \Pr[D(X) = 1 \mid f(x) = 0] \leq (1 - b) - \gamma\leq b + 2\eps -
  \gamma.
\]
  Next, we observe that
  $\min\{\Pr[f=1] , \Pr[f=0]\} \geq \gamma - 2\eps$. This follows
  since
  $\Pr[D(x,x') = f] = \Pr[D(x,x') = f = 1] + \Pr[D(x,x') = f = 0] \leq
  \Pr[f = 1] + \Pr[D = 0]$. Furthermore,
  $$
  \Pr[f = 1] \geq \Pr[D(x,x') = f] - \Pr[D = 0] \geq b + \gamma - (1-
  b) = (2b - 1) + \gamma.
  $$
  Similarly,
  $\Pr[D(x,x') = f = 1] + \Pr[D(x,x') = f = 0] \leq \Pr[f = 0] + \Pr[D
  = 1]$
  $$
  \Pr[f = 0] \geq \Pr[D(x,x') = f] - \Pr[D = 1] \geq b + \gamma - b =
  \gamma.
  $$
  This completes our proof.
\end{proof}


\begin{theorem}[Extension of
  Theorem~\ref{thm:reduction1}]\label{thm:reduction1a}
  Fix any distribution $P$ over individual data points, any set of
  group indicators $\cG$, and any classifier $D$.  Suppose that the
  (statistical parity) base rate
  $b_{SP}(D, P) \in [1/2 - \eps, 1/2 + \eps]$ for some $\eps > 0$. The
  following two relationships hold:
\begin{itemize}

\item If $D$ is
  $(\gamma - 2\eps, \gamma - 2\eps, \alpha', \beta')$-auditable for
  statistical parity fairness under distribution $P$ and group
  indicator set $\cG$, then the class $\cF$ is
  $(\gamma, \alpha'\beta' - 3\eps)$-weakly agnostically learnable under
  $\PXD$.


\item If $\cG$ is $(\alpha\beta - 3\eps, \gamma)$-weakly agnostically
  learnable under distribution $\PXD$, then $D$ is
  $(\alpha, \beta, \gamma- 3\eps, \gamma - 3\eps)$-auditable for statistical
  parity fairness under $P$ and group indicator set $\cG$.
\end{itemize}

\end{theorem}



\begin{proof}
  Suppose that the class $\cG$ satisfies
  $\min_{f\in \cG} err(f, \PXD) \leq 1/2 - \gamma$. Then by
  \Cref{acc-dis2}, there exists some $g\in \cG$ such that
  $\Pr[g(x) = 1]\geq \gamma - 2\eps$ and
  $|\Pr[D(X) = 1 \mid g(x) = 1] - b_1(D, p)| > \gamma - 2\eps$.  By
  the assumption of auditability, we can then use the auditing
  algorithm to find a group $g'\in \cG$ that is an
  $(\alpha', \beta')$-unfair certificate of $D$.  By \Cref{dis-acc2},
  we know that either $g'$ or $\neg g'$ predicts $D$ with an accuracy
  of at least $1/2 + \alpha'\beta' - 3\eps$.


  In the reverse direction: consider the auditing problem on the
  classifier $D$.  We can treat each pair $(x, D(X))$ as a labelled
  example and learn a hypothesis in $\cG$ that approximates the
  decisions made by $D$.  Suppose that $D$ is
  $(\alpha, \beta)$-unfair. Then by \Cref{dis-acc2}, we know that
  there exists some $g\in \cG$ such that
  $\Pr[D(X) = g(x)] \geq 1/2 + \alpha\beta - 3\eps$. Therefore, the
  weak agnostic learning algorithm from the hypothesis of the theorem
  will return some $g'$ with $\Pr[D(X) = g'(x)] \geq 1/2 + \gamma$. By
  \Cref{acc-dis2}, we know $g'$ or $\neg g'$ is a
  $(\gamma - 3\eps, \gamma - 3\eps)$-unfair certificate for $D$.
\end{proof}
\fi

\iffalse
\section{Fair Learning with Statistical Parity Fairness}


In this section we show that straight forward variants of the two algorithms in
\Cref{sec:learning} also solve the learning problem for the constraint of statistical parity (SP)
fairness. The crux of both algorithms is the reduction from the best response problems for both the Learner and Auditor
to CSC problems. Then if we assume CSC oracles for $\mathcal{G}$ and $\mathcal{H}$, we can implement FTPL for the Learner and best response for the auditor, recovering \Cref{alg:fairnr} and its poly-time convergence guarantees for SP fairness. \Cref{alg:fairfict}, fictitious play, is even simpler: we simply use the CSC oracles to have both players best respond to the empirical average of the other player's strategies. We now derive the CSC reductions, starting from rewriting the constrained optimization problem for statistical parity.

\iffalse
Similar to \Cref{alg:fairfict}, the algorithm iteratively simulates a
sequence of best responses by the Learner and the Auditor. However,
since we consider SP fairness constraint, the best responses are
different.
\fi

\paragraph{Rewriting the SP Fair ERM LP.}
Analogously for statistical parity fairness we define $G_{\alpha} = \{g \in \cG | \Pr_{P}[g(x) = 1] \geq \alpha \}$. Let $p$ denote a probability distribution over
$\cH$. We aim to solve the following ``Fair ERM'' problem:
\begin{align}
  &  \min_{p\in \Delta_\cH}\; \Ex{h\sim p}{err(h, P)}\\
  \mbox{such that } \forall g\in \cG_\alpha \quad &  \Ex{h\sim p}{h(x,x')|g(x) =1} - \Ex{h\sim p}{h(x,x')} \leq \beta \label{fair1}\\
&  \Ex{h\sim p}{h(x,x')} - \Ex{h\sim p}{h(x,x')|g(x) =1} \leq \beta \label{fair2}
\end{align}
where $err(h, P) = \Pr_P[h(x, x') \neq y]$.



As before we we will write $\cG(S)$ and $\cH(S)$
to denote the set of all labellings on $S$ that are induced by $\cG_\alpha$
and $\cH$ respectively, that is
\[
  \cG(S) = \{(g(x_1), \ldots , g(x_n)) \mid g\in \cG_\alpha\} \qquad
  \mbox{and,} \qquad \cH(S) = \{(h(X_1), \ldots , h(X_n))\mid h\in
  \cH\}
\]

We now write the ERM LP as an optimization where the distribution is
over labellings in $\cH(S)$ and the set of subgroups is defined by labellings
in $\cG(S)$. We also rewrite the constraints again in a way that will reduce
the problem of finding the most violated constraint (Auditor best response) to
a cost-sensitive classification problem. We define:
\begin{itemize}
\item $\Phi_+(h,g) \equiv (\Pr[h(X) =1]-\beta)\Pr[g(x) =1]-\Pr[h(X)=1,g(x)=1]$
\item $\Phi_-(h,g) \equiv  \Pr[h(X) = 1, g(x) = 1] - \left(\Pr[h(X)=1] + \beta\right)\Pr[g(x) = 1]$
\end{itemize}

As before we now focus on the following optimization problem which is easily seen to be equivalent to the original LP:
\begin{align}
    \min_{p\in \Delta_{\cH(S)}} & \; \Ex{h\sim p}{err(h, P)}   \\ \mbox{such that for each }  g\in \cG(S):\qquad
  & \Phi_+(p, g)\leq 0 \label{sp1}\\
  & \Phi_-(p,g) \leq 0 \label{sp2}
\end{align}
We then form the partial Lagrangian of the LP, and view it as a zero-sum two player game with payoff function
\[
  U(h, \lambda) = err(h, P) + \sum_{g\in \cG(S)}
  \left(\lambda_g^+\Phi_+(h, g) + \lambda_g^-\Phi_-(h, g) \right)
\]

We now reduce the best response problems for the
primal (Learner) and dual (Auditor) players to cost-sensitive
classification problems.

\paragraph{Learner's best response.} Fix any mixed strategy (dual
solution) $\lambda$ of the Auditor. Then the Learner's best response is given by:
\[\label{eq:br2} \min_{h\in \cH} \; err(h, P) + \sum_{g\in \cG(S)}
  \left(\lambda_g^+\Phi_+(h, g) + \lambda_g^-\Phi_-(h, g) \right)
\]




We can reduce this problem to one that can be solved with a single
call to a {cost-sensitive classification oracle}. In particular, we
can find a best response for the Learner by making a call to the
cost-sensitive classification oracle, in which we assign costs to each
example $(X_i, y_i)$ as follows:
\begin{itemize}
% \item if $y_i = 1$, then $c_i^0 = \frac{1}{n}$ and $c_i^1 = 0$;
\item if $y_i = 1$, then $c_i^0 = 0$ and
  $c_i^1 = - \frac{1}{n} + \frac{1}{n} \sum_{g\in
    \cG_\alpha}({\lambda}_g^+-{\lambda}_g^-)(\Pr[g(x) =1] -
  1)\textbf{1}[g(x_i=1)]$;
\item otherwise, $c_i^0 = 0$ and
  \begin{align}
    c_i^1 = \frac{1}{n} + \frac{1}{n} \sum_{g\in \cG_\alpha}({\lambda}_g^+-{\lambda}_g^-)(\Pr[g(x) =1] - 1)\textbf{1}[g(x_i=1)]
  \end{align}
\end{itemize}
\iffalse
\begin{itemize}
\item $c_i^0 = 0$ and
  $c_i^1 = \frac{1}{n} + \frac{1}{n} \sum_{g\in \cG_\alpha}({\lambda}_g^+-{\lambda}_g^-)(\Pr[g(x) =1])\textbf{1}[g(x_i=1)]$
\end{itemize}
\fi
Crucially, we can now define $\text{LC}(\lambda)_i = c_i^{1}$, and write the Learner's optimization problem as:
$$
 \min_{p\in \Delta_{\cH(S)}} \langle p, \text{LC}(\lambda) \rangle
$$
If we fix $\lambda = \bar{\lambda}_t$, the empirical average dual vector, we can either solve the linear optimization directly using the CSC oracle (\Cref{alg:fairfict}) or solve it noisily using FTPL (\Cref{alg:fairnr}).
\paragraph{Auditor's best response.}
Fix any mixed strategy (primal solution) $p \in \Delta_{\cH(S)}$ of
the Learner. The Auditor's best response is given by:
\begin{equation} \label{eq:abr} \argmax_{\lambda \in \Lambda} \;
  err(p, P) + \sum_{g\in \cG(S)} \left( \lambda_g^+ \Phi_+(p, g) +
    \lambda_g^- \Phi_-(p, g) \right) = \argmax_{\lambda \in \Lambda}
  \sum_{g\in \cG(S)} \left( \lambda_g^+ \Phi_+(p, g) + \lambda_g^-
    \Phi_-(p, g) \right)
\end{equation}

To solve this best response problem, consider the problem of computing
$(\hat g, \hat \bullet) = \argmax_{(g, \bullet)}\Phi_\bullet(p, g)$.
There are two cases. In the first case, $p$ is a strictly feasible primal
solution: that is $\Phi_{\hat \bullet} (p, \hat g) < 0$. In this case, the
solution sets $\lambda =
\mathbf{0}$. Otherwise, if $p$ is not strictly feasible, the best response is to set
$\lambda_{\hat g}^{\hat \bullet} = C$ (and all other coordinates to 0).

Therefore, it suffices to solve for
$\argmax_{(g, \bullet)}\Phi_\bullet(p, g)$. We will proceed by solving
$\argmax_{g} \Phi_+(p, g)$ and $\argmax_{g} \Phi_-(p, g)$ separately:
 both problems can be reduced to a cost-sensitive classification
problem. To solve for $\argmax_{g} \Phi_+(p, g)$ with a CSC oracle, we
assign costs to each example $(X_i, y_i)$ as follows:
\begin{itemize}
\item  $c_i^0 = 0$ and
  \begin{align}
    c_i^1 = \frac{-1}{n}\left[ \left(\Ex{h\sim p}{h(X)} - \beta\right) - \textbf{1}\{h(X_i) = 1\}  \right]
  \end{align}
\end{itemize}

To solve for $\argmax_{g} \Phi_-(p, g)$ with a CSC oracle, we set $c_i^0 = 0$, and set
\[
  c_i^1 = \frac{-1}{n}\left[ \left(\Ex{h\sim p}{h(X_i)} + \beta \right) - \textbf{1}\{h(X_i) = 1\}\right]
\]

Now since we have reduced computing the best response for both the
Learner and Auditor to solving a CSC problem, with oracles for solving
CSC in hand (e.g. agnostic learners), it is straightforward to implement the natural variants of \Cref{alg:fairnr} and \Cref{alg:fairfict} for statistical parity fairness.
\fi



\section{Missing Proofs in \Cref{sec:learning}}

\approxmm*

\begin{proof}[Proof of \Cref{thm:approxmm}]
  Let $D^*$ be the optimal feasible solution for our constrained
  optimization problem.  Since $D^*$ is feasible, we know that
  $\cL(D^*, \hat \lambda) \leq err(D^*, \cP)$.


  We will first focus on the case where $\hat D$ is not a feasible
  solution, that is
  $$\max_{(g, \bullet) \in \cG(S) \times \{\pm\}} \Phi_\bullet(\hat D, g) > 0$$ Let
  $(\hat g, \hat \bullet)\in \argmax_{(g, \bullet)} \Phi_\bullet(\hat
  D, g)$ and let $\lambda'\in \Lambda$ be a vector with
  $(\lambda')_{\hat g}^{\hat \bullet} = C$ and all other coordinates
  zero. By \Cref{lem:dualbr}, we know that
  $\lambda' \in \argmax_{\lambda\in \Lambda} \cL(\hat D, \lambda)$.
  \iffalse Let
  $\hat g\in \arg\max_{g\in \cG(S)} |\Ex{h\sim \hat D}{\FP(h, g)} -
  \Ex{h\sim \hat D}{\FP(h)}|$, and let
  $\lambda'\in \RR^{|\cG_\alpha|}$ be the best response to $\hat D$
  with
  $$\lambda'_{\hat g} = C \, \, \text{sgn}\left(\Ex{h\sim \hat D}{\FP(h, \hat g)}
    - \Ex{h\sim \hat D}{\FP(h)}\right)$$ and zeros in all other
  coordinates.  % \ar{We should note somewhere in a lemma that the best response for the dual player puts all its weight on one group.}
  \fi By the definition of a $\nu$-approximate minmax solution, we
  know that
  $\cL(\hat D, \hat \lambda) \geq \cL(\hat D, \lambda') - \nu$. This
  implies that
  \begin{equation}\label{batman}
    \cL(\hat D, \hat \lambda) \geq err(\hat D, \cP) + C \, \Phi_{\hat \bullet}(\hat D, \hat g) -
    \nu
  \end{equation}
  Note that $\cL(D^*, \hat \lambda) \leq err(D^*, \cP)$, and so
  \begin{equation}\label{superman}
    \cL(\hat D, \hat \lambda) \leq \min_{D\in \Delta_{\cH(S)}}\cL(D, \hat
    \lambda) + \nu \leq \cL(D^*, \hat
    \lambda) + \nu
\end{equation}
Combining \Cref{batman,superman}, we get
  \[
    err(\hat D, \cP) + C\, \Phi_{\hat \bullet}(\hat D, \hat g) \leq
    \cL(\hat D, \hat \lambda) + \nu \leq \cL(D^*, \hat \lambda) + 2\nu
    \leq err(D^*, \cP) + 2\nu
  \]
  Note that $C\, \Phi_{\hat \bullet}(\hat D, \hat g) \geq 0$, so we
  must have $err(\hat D, \cP) \leq err(D^*, \cP) +2 \nu = \OPT + 2\nu$.
  Furthermore, since $err(\hat D, \cP), err(D^*, \cP)\in [0, 1]$, we know
  \[
    C\, \Phi_{\hat \bullet}(\hat D, \hat g) \leq 1 + 2\nu,
  \]
  which implies that maximum constraint violation satisfies
  $\Phi_{\hat \bullet}(\hat D, \hat g) \leq (1+2\nu)/C$.
  By applying \Cref{cl:stuff}, we get
  \[
    \alpha_{FP}(g, \cP)\; \beta_{FP}(g, \hat D, \cP) \leq \gamma + \frac{1 + 2\nu}{C}.
  \]


  Now let us consider the case in which $\hat D$ is a feasible
  solution for the optimization problem. Then it follows that there is
  no constraint violation by $\hat D$ and
  $\max_\lambda \cL(\hat D, \lambda) = err(\hat D, \cP)$, and so
  \[
    err(\hat D, \cP) = \max_\lambda \cL(\hat D, \lambda) \leq \cL(\hat
    D, \hat \lambda) + \nu \leq \min_{D} \cL(D, \hat \lambda) + 2\nu
    \leq \cL(D^*, \hat \lambda) + 2\nu \leq err(D^*, \cP) + 2\nu
  \]
  Therefore, the stated bounds hold for both cases.
\end{proof}


\dualbr*

\begin{proof}[Proof of \Cref{lem:dualbr}]
Observe:
\begin{align*}
  \argmax_{\lambda\in \Lambda} \cL(\overline D, \lambda) &=
                                                           \argmax_{\lambda\in \Lambda} \Ex{h\sim \overline
                                                           D}{err(h, \cP)} + \sum_{g\in \cG(S)} \left( \lambda_g^+\,  \Phi_+(\overline D, g) +  \lambda_g^- \,  \Phi_-(\overline D,g)\right)\\
                                                         &=   \argmax_{\lambda\in \Lambda}  \sum_{g\in \cG} \left( \lambda_g^+\,  \Phi_+(\overline D, g) +  \lambda_g^- \,  \Phi_-(\overline D,g)\right)
\end{align*}
Note that this is a linear optimization problem over the non-negative
orthant of a scaling of the $\ell_1$ ball, and so has a solution at a
vertex, which corresponds to a single group $g \in \cG(S)$. Thus,
there is always a best response $\lambda'$ that puts all the weight
$C$ on the coordinate $(\lambda')_g^\bullet$ that maximizes
$\Phi_\bullet(\overline D, g)$.\iffalse Note that
\begin{align*}
  \argmax_{\lambda\in \Lambda_{\text{pure}}} \cL(\overline D, \lambda) &=
                                                                         \argmax_{\lambda\in \Lambda_{\text{pure}}} \Ex{h\sim \overline
                                                                         D}{err(h, \cP)} + \sum_{g\in \cG(S)} \lambda_g \left(\Ex{h\sim
                                                                         \overline D} {\FP(h, g)} - \Ex{h\sim \overline D}{\FP(h)}\right)\\
                                                                       &=   \argmax_{\lambda\in \Lambda_{\text{pure}}}  \sum_{g\in \cG(S)} \lambda_g \left(\Ex{h\sim
                                                                         \overline D} {\FP(h, g)} - \Ex{h\sim \overline D}{\FP(h)}\right)
\end{align*}
Note that this is a linear optimization problem over a scaling of the
$\ell_1$ ball, and so has a solution at a vertex, which corresponds to
a single group $g \in \cG_\alpha$. Thus, there is always a best
response $\lambda'$ that puts all its weight on the group that
maximizes
$ |\Ex{h\sim \overline D}{\FP(h, g)} - \Ex{h\sim \overline
  p}{\FP(h)}|$.  \fi
\end{proof}

\learneregret*

\begin{proof}[Proof of \Cref{learner-regret}]
  To instantiate the regret bound in \Cref{thm:ftpl-regret}, we just
  need to provide a bound on the maximum absoluate value over the
  coordinates of the loss vector (the quantity $M$ in
  \Cref{thm:ftpl-regret}). For any $\lambda\in \Lambda$, the absolute
  value of the $i$-th coordinate of $\LC(\lambda)$ is bounded by:
  \begin{align*}
    &\frac{1}{n} + \left|\frac{1}{n}\sum_{g\in \cG(S)}
      (\lambda_g^+ - \lambda_g^-) \left(\Pr[g(x) = 1\mid y = 0] -  1\right)
      \, \mathbf{1}[g(x_i) = 1]\right|\\
    \leq & \frac{1}{n} + \frac{1}{n} \left(\sum_{g\in \cG(S)}\left| \lambda_g^+ - \lambda_g^- \right|\right) \max_{g\in \cG(S)} \left( \Pr[g(x) = 1\mid y=0] \mathbf{1}g(x_i)=1 \right)\\
    \leq &\frac{1}{n} + \frac{1}{n} \left(\sum_{g\in \cG(S)}\left| \lambda_g^+ \right| +\left| \lambda_g^- \right|\right) \leq \frac{1 + C}{n}
  \end{align*}
  Also note that the dimension of the optimization is the size of the
  dataset $n$. This means if we set
  $\eta = \frac{n}{(1+C)} \sqrt{\frac{1}{\sqrt{n} T}}$, the regret of
  the learner will then be bounded by $2 n^{1/4} {(1+C)} \sqrt{T}$.
\end{proof}


\samplingerr*

\begin{proof}[Proof of \Cref{lem:sampling-err}]
  % \iffalse Recall that for any distribution $p'$ over $\cH(S)$, the
  % Auditor's best response problem is
  % $$\argmax_{(g, \bullet)\in \cG(S)\times \{\pm\}} \Phi_\bullet(D',
  % g).$$ \fi
  Recall that for any distribution $D'$ over $\cH(S)$ the expected
  payoff function is defined as
  \begin{align*}
    \Ex{h\sim \hat D}{U(h, \lambda)} -  \Ex{h\sim D}{U(h, \lambda)} &= \Ex{h\sim \hat D}{err(h, \cP)} 
                                                                      + \Ex{h\sim \hat D}{\sum_{g\in \cG(S)}
                                                                      \left(\lambda_g^+\Phi_+(h, g) + \lambda_g^-\Phi_-(h, g) \right)}\\
&- \Ex{h\sim D}{err(h, \cP)} 
                                                                      + \Ex{h\sim D}{\sum_{g\in \cG(S)}
                                                                      \left(\lambda_g^+\Phi_+(h, g) + \lambda_g^-\Phi_-(h, g) \right)}
  \end{align*}
  By the triangle inequality, it suffices to show that with
  probability $(1 - \delta)$,
  $A = |\Ex{h\sim D}{err(h, \cP)} - \Ex{h\sim \hat D}{err(h, \cP)}|
  \leq \xi/2$ and for all $\lambda\in \Lambda$ and $g\in \cG(S)$,
  \[
    B = \left| \Ex{h\sim \hat D}{\sum_{g\in \cG(S)}
        \left(\lambda_g^+\Phi_+(h, g) + \lambda_g^-\Phi_-(h, g)
        \right)} - \Ex{h\sim D}{\sum_{g\in \cG(S)}
        \left(\lambda_g^+\Phi_+(h, g) + \lambda_g^-\Phi_-(h, g)
        \right)}\right| \leq \xi/2
  \]



    The first part follows directly from a simple application of the
    Chernoff-Hoeffding bound (\Cref{chernoff}): with probability
    $(1 - \delta/2)$, $A \leq \xi/2$, as long as
    $m\geq 2\ln(4/\delta)/\xi^2 $.

    To bound the second part, we first note that by H\"{o}lder's
    inequality, we have
    $$B\leq \|\lambda\|_1 \max_{(g, \bullet)\in \cG(S) \times
      \{\pm\}}|\Phi_\bullet(D, g) - \Phi_\bullet(\hat D, g)|$$ 

    Since for all $\lambda\in \Lambda$ we have $\|\lambda\|_1 \leq C$,
    it suffices to show that with probability $1 - \delta/2$,
    $|\Phi_\bullet(D, g) - \Phi_\bullet(\hat D, g)| \leq \xi/(2C)$
    holds for all $\bullet\in \{-, +\}$ and $g\in \cG(S)$. Note that
  \begin{align*}
    |\Phi_\bullet(D, g) - \Phi_\bullet(\hat D, g)| = \left|\left(\Ex{h\sim
    D}{ \FP(h)} - \Ex{h\sim \hat D}{\FP(h)} \right)\Pr[y=0 , g(x) =
    1] \right. \\+ \left. \left(\Ex{h\sim D}{\Pr[h(X) = 1, y= 0, g(x) = 1 ]} -
    \Ex{h\sim \hat D}{\Pr[h(X) = 1, y= 0, g(x) = 1 ]} \right)\right|
  \end{align*}
  We can rewrite the absolute value of first term:
  \begin{align*}
    &    \left|\left(\Ex{h\sim D}{ \FP(h)} - \Ex{h\sim \hat D}{\FP(h)}
      \right)\Pr[y=0 , g(x) = 1] \right|\\ = &\left|\left(\Ex{h\sim D}{ \Pr[h(X) = 1\mid y=0]} -
                                               \Ex{h\sim \hat D}{\Pr[h(X) = 1\mid y = 0]} \right)\Pr[g(x) =  1\mid y=0]\right|\\
    \leq  &\left|\left(\Ex{h\sim D}{ \Pr[h(X) = 1,  y=0]} -
            \Ex{h\sim \hat D}{\Pr[h(X) = 1, y = 0]} \right)\right|
  \end{align*}
  where the last inequality follows from
  $\Pr[g(x) = 1\mid y=0] \leq 1$.

  Note that
  $\Ex{h\sim \hat D}{\Pr[h(X) = 1, y= 0, g(x) = 1 ]} = \frac{1}{m}
  \sum_{j=1}^m \Pr[h^j(X) = 1, y= 0, g(x) = 1 ]$, which is an average
  of $m$ i.i.d. random variables with expectation
  $\Ex{h\sim D}{\Pr[h(X) = 1, y= 0, g(x) = 1 ]}$. By the
  Chernoff-Hoeffding bound (\Cref{chernoff}), we have
  \begin{equation}
    \Pr\left[ \left|\Ex{h\sim D}{ \Pr[h(X) = 1, y=0]} -
        \Ex{h\sim \hat D}{\Pr[h(X) = 1, y = 0]} \right| >
      \frac{\xi}{4C} \right] \leq 2\exp\left(-\frac{ \xi^2 m}{ 8 C^2}
    \right)\label{hp1}\end{equation}
  In the following, we will let $\delta_0 = 2\exp\left(-\frac{ \xi^2 m}{ 8 C^2}
  \right)$.  Similarly, we also have for each $g\in \cG(S)$,
  \begin{equation}
    \Pr\left[\left| \Ex{h\sim D}{\Pr[h(X) = 1, y= 0, g(x) = 1 ]} -
        \Ex{h\sim \hat D}{\Pr[h(X) = 1, y= 0, g(x) = 1 ]}\right| >
      \frac{\xi}{4C}\right] \leq \delta_0\label{hp2}
  \end{equation}


% large enough such that
%   $\delta/2 \geq (|\cG(S)| + 1) \delta_0$,
  By taking the union bound over \eqref{hp1} and \eqref{hp2} over all
  choices of $g\in \cG(S)$, we have with probability at least
  $(1 - \delta_0(1 + |\cG(S)|))$,
  \begin{equation} \left|\Ex{h\sim D}{ \Pr[h(X) = 1, y=0]} - \Ex{h\sim
        \hat D}{\Pr[h(X) = 1, y = 0]} \right| \leq \frac{\xi}{4C}
    \quad
\label{wonder}\end{equation}
{and,}
  \begin{equation}
    \left| \Ex{h\sim D}{\Pr[h(X) = 1, y= 0, g(x) = 1 ]} - \Ex{h\sim
        \hat D}{\Pr[h(X) = 1, y= 0, g(x) = 1 ]}\right| \leq
    \frac{\xi}{4C} \quad \mbox{for all } g \in \cG(S). \label{woman}
  \end{equation}
  Note that by Sauer's lemma (\Cref{lem:sauer}),
  $|\cG(S)|\leq O\left( n^{d_2} \right)$. Thus, there exists an
  absolute constant $c_0$ such that
  $m \geq c_0\frac{C^2\left(\ln(1/\delta) + d_2\ln(n)\right)}{\xi^2}$
  implies that failure probability above
  $\delta_0(1 + |\cG(S)|) \leq \delta/2$. We will assume $m$
  satisifies such a bound, and so the events of \eqref{wonder} and
  \eqref{woman} hold with probaility at least $(1 - \delta/2)$.  Then
  by the triangle inequality we have for all
  $(g, \bullet)\in \cG(S)\times\{\pm\}$,
  $|\Phi_\bullet(D, g) - \Phi_\bullet(\hat D, g)| \leq \xi/(2C)$,
  which implies that $B \leq \xi/2$. This completes the proof.
  %% delta = delta_0 * O(n^d_2)
\end{proof}

\iffalse
The following corollary follows directly from \Cref{lem:sampling-err}.

\begin{corollary}
  Fix any $\xi, \delta\in (0, 1)$ and any distribution $D$ over
  $\cH(S)$. Let $h^1, \ldots, h^m$ be $m$ i.i.d. draws from $p$, and
  let $\hat D$ be the empirical distribution over the realized
  sample. Let
  $$\hat{\lambda} \in \argmax_{\lambda'\in \Lambda} \Ex{h\sim \hat
    D}{U(h, \lambda')}$$ be the Auditor's best response against
  $\hat D$. Then with probability at least $1 - \delta$ over the
  random draws of $h^j$'s,
  \[
    \max_{\lambda} \Ex{h\sim D}{U(h, \lambda)} -\xi \leq \Ex{h\sim
      D}{U(h, \hat \lambda)},
  \]
  as long as
  $m \geq c_0\frac{C^2\left(\ln(1/\delta) +
      d_2\ln(n)\right)}{\xi^2}$ for some absolute constant $c_0$
  and $d_2 = \VCD(\cG)$.
\end{corollary}
\fi

\begin{claim}\label{cor:approx-br}
  Suppose there are two distributions $D$ and $\hat D$ over $\cH(S)$
  such that
  \[
    \max_{\lambda\in \Lambda} \left| \Ex{h\sim \hat D}{U(h, \lambda)}
      - \Ex{h\sim D}{U(h, \lambda)} \right| \leq \xi.
  \]
Let
  $$\hat{\lambda} \in \argmax_{\lambda'\in \Lambda} \Ex{h\sim \hat
    D}{U(h, \lambda')}$$ 
Then
  \[
    \max_{\lambda} \Ex{h\sim D}{U(h, \lambda)} -\xi \leq \Ex{h\sim
      D}{U(h, \hat \lambda)},
  \]
\end{claim}


\auditorbr*

\begin{proof}
  Let $\gamma_A^t$ be defined as
  \[
    \gamma_A^t = \max_{\lambda\in \Lambda} \left| \Ex{h\sim \hat
        D^t}{U(h, \lambda)} - \Ex{h\sim D^t}{U(h, \lambda)} \right|
  \]
  By instantiating \Cref{lem:sampling-err} and applying union bound
  across all $T$ steps, we know with probability at least
  $1 - \delta$, the following holds for all $t\in [T]$:
  \[
    {\gamma_A^t} \leq \sqrt{\frac{c_0 C^2 (\ln(T/\delta) + d_2
        \ln(n))}{m}}
  \]
  where $c_0$ is the absolute constant in \Cref{lem:sampling-err} and
  $d_2 = \VCD(\cG)$. 

  Note that by \Cref{cor:approx-br}, the Auditor is performing a
  $\gamma_A^t$-approximate best response at each round $t$. Then we
  can bound the Auditor's regret as follows:
  \begin{align*}
    \gamma_A = \frac{1}{T}\left[\max_{\lambda\in \Lambda} \sum_{t=
    1}^T \Ex{h\sim D^t}{U(h, \lambda)} - \sum_{t=1}^T \Ex{h\sim
    D^t}{U(h, \lambda^t)}\right] &\leq \frac{1}{T}\sum_{t=
                                   1}^T\left( \max_{\lambda\in \Lambda} \Ex{h\sim D^t}{U(h, \lambda)} - \Ex{h\sim
                                   D^t}{U(h, \lambda^t)} \right)  \\
                                 &\leq \max_T \gamma_A^t
  \end{align*}
  It follows that with probability $1 - \delta$, we have
\begin{align*}
  \gamma_A &\leq \sqrt{\frac{c_0 C^2 (\ln(T/\delta) + d_2
        \ln(n))}{m}}
\end{align*}
which completes the proof.
\end{proof}

\iffalse
\polytime*

\begin{proof}[Proof of \Cref{thm:polytime}]
  By \Cref{thm:approxmm}, it suffices to show that with probability at
  least $1 - \delta$, $(\hat D, \overline \lambda)$ is a
  $\nu$-approximate equilibrium in the zero-sum game. As a first step,
  we will rely on \Cref{fs} to show that
  $(\overline D, \overline \lambda)$ forms an approximate equilibrium.

  By \Cref{learner-regret}, the regret of the sequence
  $D^1, \ldots , D^T$ is bounded by:
  \[
    \gamma_L = \frac{1}{T}\left[\sum_{t= 1}^T \Ex{h\sim D^t}{U(h,
        \lambda^t)} - \min_{h\in \cH(S)}\sum_{t=1}^T U(h, \lambda^t)
    \right] \leq \frac{2 n^{1/4} {(1+C)}}{\sqrt{T}}
  \]

  % Let $\gamma_A^t$ be defined as
  % \[
  %   \gamma_A^t = \max_{\lambda\in \Lambda} \left| \Ex{h\sim \hat
  %       D^t}{U(h, \lambda)} - \Ex{h\sim D^t}{U(h, \lambda)} \right|
  % \]
  % By instantiating \Cref{lem:sampling-err} and applying union bound
  % across all $T$ steps, we know with probability at least
  % $1 - \delta/2$, the following holds for all $t\in [T]$:
  % \[
  %   {\gamma_A^t} \leq \sqrt{\frac{c_0 C^2 (\ln(2T/\delta) + d_2
  %       \ln(n))}{m}}
  % \]
  % where $c_0$ is the absolute constant in \Cref{lem:sampling-err} and
  % $d_2 = \VCD(\cG)$. 

  % Note that by \Cref{cor:approx-br}, the Auditor is performing a
  % $\gamma_A^t$-approximate best response at each round $t$. Then we
  % can bound the Auditor's regret as follows:
  % \begin{align*}
  %   \gamma_A = \frac{1}{T}\left[\max_{\lambda\in \Lambda} \sum_{t=
  %   1}^T \Ex{h\sim D^t}{U(h, \lambda)} - \sum_{t=1}^T \Ex{h\sim
  %   D^t}{U(h, \lambda^t)}\right] &\leq \frac{1}{T}\sum_{t=
  %                                  1}^T\left( \max_{\lambda\in \Lambda} \Ex{h\sim D^t}{U(h, \lambda)} - \Ex{h\sim
  %                                  D^t}{U(h, \lambda^t)} \right)  \\
  %                                &\leq \max_T \gamma_A^t
  % \end{align*}
  By \Cref{lem:auditor-br}, with probability $1 - \delta/2$, we have
\begin{align*}
  \gamma_A &\leq \sqrt{\frac{c_0 C^2 (\ln(2T/\delta) + d_2
        \ln(n))}{m}}
\end{align*}
  We will condition on this upper-bound event on $\gamma_A$ for the
  rest of this proof, which is the case except with probability
  $\delta/2$.

  By \Cref{fs}, we know that the average plays
  $(\overline D, \overline \lambda)$ form an
  $(\gamma_L+\gamma_A)$-approximate equilibrium. By
  \Cref{lem:sampling-err}, we know that with probability at least
  $(1- \delta/2)$, the pair $(\hat D, \overline \lambda)$ form a
  $R$-approximate equilibrium, with
  \begin{align*}
    R &\leq \gamma_A + \gamma_L + \frac{\sqrt{c_0 C^2(\ln(2/\delta)
        +d_2\ln(n))}}{\sqrt{m}}\\
      & \leq  \sqrt{\frac{c_0 C^2 (\ln(2T/\delta) + d_2
        \ln(n))}{m}} 
        + \frac{\sqrt{c_0 C^2(\ln(2/\delta)
        +d_2\ln(n))}}{\sqrt{m}} + \frac{2 n^{1/4} {(1+C)}}{\sqrt{T}}\\
      & \leq 
        2\sqrt{\frac{c_0 C^2 (\ln(2T/\delta) + d_2
        \ln(n))}{m}}
        + \frac{2 n^{1/4} {(1+C)} \sqrt{\ln(2/\delta)}}{\sqrt{T}}
  \end{align*}
  Note that $R\leq \nu$ as long as we have $C = 1/\nu$,
  \[
    m = \frac{\left(\ln(2T/\delta)d_2\ln(n)\right) C^2 c_0 T}{\sqrt{n} (1+C)^2 \ln(2/\delta)}
 \quad
    \mbox{and, }\quad T = \frac{4\sqrt{n} \ln(2/\delta)}{\nu^4}
  \]
 This completes our proof.
\end{proof}

\fi




\iffalse
\section{Necessary Conditions on the Base Rates}
\sw{add the example here}
\fi


%%% Local Variables:
%%% mode: latex
%%% TeX-master: "main.tex"
%%% End:
